\IfFileExists{stacks-project.cls}{%
\documentclass{stacks-project}
}{%
\documentclass{amsart}
}

% For dealing with references we use the comment environment
\usepackage{verbatim}
\newenvironment{reference}{\comment}{\endcomment}
%\newenvironment{reference}{}{}
\newenvironment{slogan}{\comment}{\endcomment}
\newenvironment{history}{\comment}{\endcomment}

% For commutative diagrams we use Xy-pic
\usepackage[all]{xy}

% We use 2cell for 2-commutative diagrams.
\xyoption{2cell}
\UseAllTwocells

% We use multicol for the list of chapters between chapters
\usepackage{multicol}

% This is generally recommended for better output
\usepackage{lmodern}
\usepackage[T1]{fontenc}

% For cross-file-references
\usepackage{xr-hyper}

% Package for hypertext links:
\usepackage{hyperref}

% For any local file, say "hello.tex" you want to link to please
% use \externaldocument[hello-]{hello}
\externaldocument[introduction-]{introduction}
\externaldocument[conventions-]{conventions}
\externaldocument[sets-]{sets}
\externaldocument[categories-]{categories}
\externaldocument[topology-]{topology}
\externaldocument[sheaves-]{sheaves}
\externaldocument[sites-]{sites}
\externaldocument[stacks-]{stacks}
\externaldocument[fields-]{fields}
\externaldocument[algebra-]{algebra}
\externaldocument[brauer-]{brauer}
\externaldocument[homology-]{homology}
\externaldocument[derived-]{derived}
\externaldocument[simplicial-]{simplicial}
\externaldocument[more-algebra-]{more-algebra}
\externaldocument[smoothing-]{smoothing}
\externaldocument[modules-]{modules}
\externaldocument[sites-modules-]{sites-modules}
\externaldocument[injectives-]{injectives}
\externaldocument[cohomology-]{cohomology}
\externaldocument[sites-cohomology-]{sites-cohomology}
\externaldocument[dga-]{dga}
\externaldocument[dpa-]{dpa}
\externaldocument[sdga-]{sdga}
\externaldocument[hypercovering-]{hypercovering}
\externaldocument[schemes-]{schemes}
\externaldocument[constructions-]{constructions}
\externaldocument[properties-]{properties}
\externaldocument[morphisms-]{morphisms}
\externaldocument[coherent-]{coherent}
\externaldocument[divisors-]{divisors}
\externaldocument[limits-]{limits}
\externaldocument[varieties-]{varieties}
\externaldocument[topologies-]{topologies}
\externaldocument[descent-]{descent}
\externaldocument[perfect-]{perfect}
\externaldocument[more-morphisms-]{more-morphisms}
\externaldocument[flat-]{flat}
\externaldocument[groupoids-]{groupoids}
\externaldocument[more-groupoids-]{more-groupoids}
\externaldocument[etale-]{etale}
\externaldocument[chow-]{chow}
\externaldocument[intersection-]{intersection}
\externaldocument[pic-]{pic}
\externaldocument[weil-]{weil}
\externaldocument[adequate-]{adequate}
\externaldocument[dualizing-]{dualizing}
\externaldocument[duality-]{duality}
\externaldocument[discriminant-]{discriminant}
\externaldocument[derham-]{derham}
\externaldocument[local-cohomology-]{local-cohomology}
\externaldocument[algebraization-]{algebraization}
\externaldocument[curves-]{curves}
\externaldocument[resolve-]{resolve}
\externaldocument[models-]{models}
\externaldocument[functors-]{functors}
\externaldocument[equiv-]{equiv}
\externaldocument[pione-]{pione}
\externaldocument[etale-cohomology-]{etale-cohomology}
\externaldocument[proetale-]{proetale}
\externaldocument[relative-cycles-]{relative-cycles}
\externaldocument[more-etale-]{more-etale}
\externaldocument[trace-]{trace}
\externaldocument[crystalline-]{crystalline}
\externaldocument[spaces-]{spaces}
\externaldocument[spaces-properties-]{spaces-properties}
\externaldocument[spaces-morphisms-]{spaces-morphisms}
\externaldocument[decent-spaces-]{decent-spaces}
\externaldocument[spaces-cohomology-]{spaces-cohomology}
\externaldocument[spaces-limits-]{spaces-limits}
\externaldocument[spaces-divisors-]{spaces-divisors}
\externaldocument[spaces-over-fields-]{spaces-over-fields}
\externaldocument[spaces-topologies-]{spaces-topologies}
\externaldocument[spaces-descent-]{spaces-descent}
\externaldocument[spaces-perfect-]{spaces-perfect}
\externaldocument[spaces-more-morphisms-]{spaces-more-morphisms}
\externaldocument[spaces-flat-]{spaces-flat}
\externaldocument[spaces-groupoids-]{spaces-groupoids}
\externaldocument[spaces-more-groupoids-]{spaces-more-groupoids}
\externaldocument[bootstrap-]{bootstrap}
\externaldocument[spaces-pushouts-]{spaces-pushouts}
\externaldocument[spaces-chow-]{spaces-chow}
\externaldocument[groupoids-quotients-]{groupoids-quotients}
\externaldocument[spaces-more-cohomology-]{spaces-more-cohomology}
\externaldocument[spaces-simplicial-]{spaces-simplicial}
\externaldocument[spaces-duality-]{spaces-duality}
\externaldocument[formal-spaces-]{formal-spaces}
\externaldocument[restricted-]{restricted}
\externaldocument[spaces-resolve-]{spaces-resolve}
\externaldocument[formal-defos-]{formal-defos}
\externaldocument[defos-]{defos}
\externaldocument[cotangent-]{cotangent}
\externaldocument[examples-defos-]{examples-defos}
\externaldocument[algebraic-]{algebraic}
\externaldocument[examples-stacks-]{examples-stacks}
\externaldocument[stacks-sheaves-]{stacks-sheaves}
\externaldocument[criteria-]{criteria}
\externaldocument[artin-]{artin}
\externaldocument[quot-]{quot}
\externaldocument[stacks-properties-]{stacks-properties}
\externaldocument[stacks-morphisms-]{stacks-morphisms}
\externaldocument[stacks-limits-]{stacks-limits}
\externaldocument[stacks-cohomology-]{stacks-cohomology}
\externaldocument[stacks-perfect-]{stacks-perfect}
\externaldocument[stacks-introduction-]{stacks-introduction}
\externaldocument[stacks-more-morphisms-]{stacks-more-morphisms}
\externaldocument[stacks-geometry-]{stacks-geometry}
\externaldocument[moduli-]{moduli}
\externaldocument[moduli-curves-]{moduli-curves}
\externaldocument[examples-]{examples}
\externaldocument[exercises-]{exercises}
\externaldocument[guide-]{guide}
\externaldocument[desirables-]{desirables}
\externaldocument[coding-]{coding}
\externaldocument[obsolete-]{obsolete}
\externaldocument[fdl-]{fdl}
\externaldocument[index-]{index}

% Theorem environments.
%
\theoremstyle{plain}
\newtheorem{theorem}[subsection]{Theorem}
\newtheorem{proposition}[subsection]{Proposition}
\newtheorem{lemma}[subsection]{Lemma}

\theoremstyle{definition}
\newtheorem{definition}[subsection]{Definition}
\newtheorem{example}[subsection]{Example}
\newtheorem{exercise}[subsection]{Exercise}
\newtheorem{situation}[subsection]{Situation}

\theoremstyle{remark}
\newtheorem{remark}[subsection]{Remark}
\newtheorem{remarks}[subsection]{Remarks}

\numberwithin{equation}{subsection}

% Macros
%
\def\lim{\mathop{\mathrm{lim}}\nolimits}
\def\colim{\mathop{\mathrm{colim}}\nolimits}
\def\Spec{\mathop{\mathrm{Spec}}}
\def\Hom{\mathop{\mathrm{Hom}}\nolimits}
\def\Ext{\mathop{\mathrm{Ext}}\nolimits}
\def\SheafHom{\mathop{\mathcal{H}\!\mathit{om}}\nolimits}
\def\SheafExt{\mathop{\mathcal{E}\!\mathit{xt}}\nolimits}
\def\Sch{\mathit{Sch}}
\def\Mor{\mathop{\mathrm{Mor}}\nolimits}
\def\Ob{\mathop{\mathrm{Ob}}\nolimits}
\def\Sh{\mathop{\mathit{Sh}}\nolimits}
\def\NL{\mathop{N\!L}\nolimits}
\def\CH{\mathop{\mathrm{CH}}\nolimits}
\def\proetale{{pro\text{-}\acute{e}tale}}
\def\etale{{\acute{e}tale}}
\def\QCoh{\mathit{QCoh}}
\def\Ker{\mathop{\mathrm{Ker}}}
\def\Im{\mathop{\mathrm{Im}}}
\def\Coker{\mathop{\mathrm{Coker}}}
\def\Coim{\mathop{\mathrm{Coim}}}

% Boxtimes
%
\DeclareMathSymbol{\boxtimes}{\mathbin}{AMSa}{"02}

%
% Macros for moduli stacks/spaces
%
\def\QCohstack{\mathcal{QC}\!\mathit{oh}}
\def\Cohstack{\mathcal{C}\!\mathit{oh}}
\def\Spacesstack{\mathcal{S}\!\mathit{paces}}
\def\Quotfunctor{\mathrm{Quot}}
\def\Hilbfunctor{\mathrm{Hilb}}
\def\Curvesstack{\mathcal{C}\!\mathit{urves}}
\def\Polarizedstack{\mathcal{P}\!\mathit{olarized}}
\def\Complexesstack{\mathcal{C}\!\mathit{omplexes}}
% \Pic is the operator that assigns to X its picard group, usage \Pic(X)
% \Picardstack_{X/B} denotes the Picard stack of X over B
% \Picardfunctor_{X/B} denotes the Picard functor of X over B
\def\Pic{\mathop{\mathrm{Pic}}\nolimits}
\def\Picardstack{\mathcal{P}\!\mathit{ic}}
\def\Picardfunctor{\mathrm{Pic}}
\def\Deformationcategory{\mathcal{D}\!\mathit{ef}}
\usepackage{cleveref}
\crefname{originallemma}{original lemma}{original lemmas}
\Crefname{originallemma}{Original lemma}{Original lemmas}
\makeatletter
\newlabel{lemma-stalk-abelian-presheaf}{{007F}{source}{}{sourcepointer-1}{}}
\newlabel{lemma-stalk-abelian-presheaf@cref}{{[originallemma][0][]007F}{[1][0][]source}{}{}{}}
\makeatother
\crefname{originalremark}{original remark}{original remarks}
\Crefname{originalremark}{Original remark}{Original remarks}
\makeatletter
\newlabel{example-sheaf-product-pointwise-algebraic-structure}{{007R}{source}{}{sourcepointer-2}{}}
\newlabel{example-sheaf-product-pointwise-algebraic-structure@cref}{{[originalremark][0][]007R}{[1][0][]source}{}{}{}}
\makeatother
\crefname{originalsection}{original section}{original sections}
\Crefname{originalsection}{Original section}{Original sections}
\makeatletter
\newlabel{section-sheafification}{{007X}{source}{}{sourcepointer-3}{}}
\newlabel{section-sheafification@cref}{{[originalsection][0][]007X}{[1][0][]source}{}{}{}}
\makeatother
\crefname{originallemma}{original lemma}{original lemmas}
\Crefname{originallemma}{Original lemma}{Original lemmas}
\makeatletter
\newlabel{lemma-diagram-fibre-product}{{0084}{source}{}{sourcepointer-4}{}}
\newlabel{lemma-diagram-fibre-product@cref}{{[originallemma][0][]0084}{[1][0][]source}{}{}{}}
\makeatother
\crefname{originallemma}{original lemma}{original lemmas}
\Crefname{originallemma}{Original lemma}{Original lemmas}
\makeatletter
\newlabel{lemma-stalk-sheafification}{{007Z}{source}{}{sourcepointer-5}{}}
\newlabel{lemma-stalk-sheafification@cref}{{[originallemma][0][]007Z}{[1][0][]source}{}{}{}}
\makeatother
\crefname{originallemma}{original lemma}{original lemmas}
\Crefname{originallemma}{Original lemma}{Original lemmas}
\makeatletter
\newlabel{lemma-check-homomorphism-stalks}{{007W}{source}{}{sourcepointer-6}{}}
\newlabel{lemma-check-homomorphism-stalks@cref}{{[originallemma][0][]007W}{[1][0][]source}{}{}{}}
\makeatother
\title{Sheafification of an abelian presheaf carries a unique abelian sheaf structure and is left adjoint}
\author{The Stacks Project Authors}
\date{stacks-project:sheaves.tex:local-snapshot-20261012; GNU FDL 1.2 or later}
\begin{document}
\maketitle
\setcounter{subsection}{1}
\noindent\textit{Editorial note.} All delivered mathematical text is the author's own, copied from the pinned Stacks Project chapter file sheaves.tex; no statement or proof step was rewritten, reordered or removed. The counted span is the author environment and its own proof environment. References inside the span are handed to the assembler as navigation stubs that carry the canonical tag of the cited tag only; whether each of them is load-bearing is recorded on the stub itself.\par\medskip
\noindent\textit{Allowed standard prerequisites (editorial).}\begin{itemize}
\item Lemma lemma-stalk-abelian-presheaf (canonical tag 007F): stalks of abelian presheaves are abelian groups and the stalk functor is compatible with the algebraic structure. cited at proof lines 1955-1957 and 1977-1979. Gives the abelian group structure on stalks used throughout.
\item Example example-sheaf-product-pointwise-algebraic-structure (canonical tag 007R): pointwise algebraic structures on products of stalks give a sheaf of algebraic structures. cited at proof lines 1957-1959. Makes Pi(F) a sheaf of abelian groups.
\item Lemma lemma-diagram-fibre-product (canonical tag 0084): the sheafification value F\textasciicircum{}\#(U) is the fibre product of a diagram of the stalks. cited at proof lines 1969-1972. Identifies F\textasciicircum{}\#(U) as a fibre product of abelian groups, hence a subgroup of Pi(F)(U).
\item Lemma lemma-stalk-sheafification (canonical tag 007Z): the stalk of the sheafification is canonically identified with the stalk of the presheaf. cited at proof lines 1979-1981 and 1989-1990. Identifies F\_x with F\textasciicircum{}\#\_x as abelian groups and lets the abelian property be checked on stalks.
\item Lemma lemma-check-homomorphism-stalks (canonical tag 007W): a map of sheaves is a homomorphism of algebraic structures iff it is on all stalks. cited at proof lines 1990-1991. Upgrades the set-level adjoint morphism to a morphism of abelian sheaves.
\end{itemize}\medskip
\ifdefined\hypertarget\hypertarget{sourcepointer-1}{}\fi
\noindent\textit{Source reference.} Original lemma 007F: sources/stacks-chow/tags.txt, source lines 279--279. Cited at proof lines 1955-1957 and 1977-1979; load-bearing inside the counted span.\par\medskip
\ifdefined\hypertarget\hypertarget{sourcepointer-2}{}\fi
\noindent\textit{Source reference.} Original remark 007R: sources/stacks-chow/tags.txt, source lines 291--291. Cited at proof lines 1957-1959 to make Pi(F) a sheaf of abelian groups; load-bearing inside the counted span (Example cited as kind=remark per kind mapping).\par\medskip
\ifdefined\hypertarget\hypertarget{sourcepointer-3}{}\fi
\noindent\textit{Source reference.} Original section 007X: sources/stacks-chow/tags.txt, source lines 297--297. Cited at proof lines 1953-1954 and 1962-1963 for the construction being used; navigation for the setting.\par\medskip
\ifdefined\hypertarget\hypertarget{sourcepointer-4}{}\fi
\noindent\textit{Source reference.} Original lemma 0084: sources/stacks-chow/tags.txt, source lines 304--304. Cited at proof lines 1969-1972 for the subgroup argument; load-bearing inside the counted span.\par\medskip
\ifdefined\hypertarget\hypertarget{sourcepointer-5}{}\fi
\noindent\textit{Source reference.} Original lemma 007Z: sources/stacks-chow/tags.txt, source lines 299--299. Cited at proof lines 1979-1981 and 1989-1990; load-bearing inside the counted span.\par\medskip
\ifdefined\hypertarget\hypertarget{sourcepointer-6}{}\fi
\noindent\textit{Source reference.} Original lemma 007W: sources/stacks-chow/tags.txt, source lines 296--296. Cited at proof lines 1990-1991 to upgrade the adjoint to abelian sheaves; load-bearing inside the counted span.\par\medskip
\begin{lemma}
\label{lemma-sheafify-abelian-presheaf}
Let $X$ be a topological space.
Let $\mathcal{F}$ be an abelian presheaf on $X$.
Then there exists a unique structure of
abelian sheaf on $\mathcal{F}^\#$ such that
$\mathcal{F} \to \mathcal{F}^\#$ is a morphism
of abelian presheaves. Moreover, the following adjointness
property holds
$$
\Mor_{\textit{PAb}(X)}(\mathcal{F}, i(\mathcal{G}))
=
\Mor_{\textit{Ab}(X)}(\mathcal{F}^\#, \mathcal{G}).
$$
\end{lemma}

\begin{proof}
Recall the sheaf of sets $\Pi(\mathcal{F})$ defined in
Section \ref{section-sheafification}. All the stalks
$\mathcal{F}_x$ are abelian groups, see
Lemma \ref{lemma-stalk-abelian-presheaf}.
Hence $\Pi(\mathcal{F})$ is a sheaf of abelian groups by
Example \ref{example-sheaf-product-pointwise-algebraic-structure}.
Also, it is clear that the map $\mathcal{F} \to \Pi(\mathcal{F})$
is a morphism of abelian presheaves. If we show that
condition $(*)$ of Section \ref{section-sheafification} defines a subgroup
of $\Pi(\mathcal{F})(U)$ for all open subsets $U \subset X$,
then $\mathcal{F}^\#$ canonically inherits the structure of abelian sheaf.
This is quite easy to do by hand, and we leave it to the
reader to find a good simple argument. The argument we use here,
which generalizes to presheaves of algebraic structures is the following:
Lemma \ref{lemma-diagram-fibre-product} show that
$\mathcal{F}^\#(U)$ is the fibre product of a diagram
of abelian groups. Thus $\mathcal{F}^\#$ is an abelian
subgroup as desired.

\medskip\noindent
Note that at this point $\mathcal{F}^\#_x$ is an abelian
group by Lemma \ref{lemma-stalk-abelian-presheaf}
and that $\mathcal{F}_x \to \mathcal{F}^\#_x$ is a
bijection (Lemma \ref{lemma-stalk-sheafification})
and a homomorphism of abelian groups. Hence
$\mathcal{F}_x \to \mathcal{F}^\#_x$ is an isomorphism
of abelian groups. This will be used below without further mention.

\medskip\noindent
To prove the adjointness property we use the adjointness
property of sheafification of presheaves of sets. For example
if $\psi : \mathcal{F} \to i(\mathcal{G})$ is morphism of presheaves
then we obtain a morphism of sheaves
$\psi' : \mathcal{F}^\# \to \mathcal{G}$. What we have to do is to check
that this is a morphism of abelian sheaves. We may do this
for example by noting that it is true on stalks,
by Lemma \ref{lemma-stalk-sheafification}, and then using
Lemma \ref{lemma-check-homomorphism-stalks} above.
\end{proof}

\end{document}
